KernelBench v0.1 and the original KernelBench paper established a useful framing: models are given a PyTorch reference implementation and must produce a faster CUDA kernel~\cite{kernelbench,stanford-kernelbenchv01}. Subsequent investigations revealed several reward-hacking patterns, including calling high-level libraries instead of writing kernels, manipulating timing, or exploiting numerically trivial tasks~\cite{metr2025}. These findings motivated a ground-up redesign rather than incremental fixes.

KernelBench v3 is designed with two principles. First, tasks should reflect modern kernel engineering needs, including fused matmul chains, attention blocks, and emerging linear-attention variants. Second, baselines should be meaningful: for some Level 4 tasks, the baseline is already an optimized Triton implementation, so beating the baseline requires genuinely competitive kernel engineering.

This redesign also emphasizes cost realism. Running large model sweeps across many problems and hardware targets can become prohibitively expensive. KernelBench v3 focuses on two widely deployed data center GPUs and a smaller, curated task set while preserving the evaluation challenge.
